\subsection{The actual nonstandard generators as geometric quasi-endomorphisms}

\begin{proposition}[P-F70: a fiberwise source-block correspondence realization]
\label{prop:nonstandardcorrespondence}
Fix $q=3^e$ and retain the ordinary abelian surface
$\mathcal A/\mathbb F_Q$, $Q=q^{116}$, from P-F69.
Write
\[
 L=\operatorname{End}^0_{\mathbb F_Q}(\mathcal A)
      =\mathbb Q(\Pi),\qquad
 K=\mathbb Q(\sqrt{A(q)})\subset L,
\]
where $\Pi$ is its $Q$-power Frobenius endomorphism.
Set
\[
 \mathsf V=Q\Pi^{-1},\qquad
 a=\frac{\Pi+\mathsf V}{2q^{26}}\in K,\qquad
 \gamma=S-\frac aF.                                           \tag{P-F70a}
\]
All source Laurent scalars $S,F,f$ in this section are evaluated
at $q$ with exactly the normalization of P-F52.
The two real conjugates of $a$ are the actual values
$\alpha_\pm(q)$, and the two conjugates of $\gamma$ are
$g_1,g_2$ from (P-F52f1).

Let $V_q$ be the rational specialization at $q$ of the free
rank-four source lattice $\mathcal V$ in (P-F52ax), and let
\[
 B_q^{\mathrm{im}}=\mathbb Q[Q_1,Q_2]
                       \subset\operatorname{End}_{\mathbb Q}(V_q)
                                                               \tag{P-F70b}
\]
be its unital image algebra.  Then
\[
 B_q^{\mathrm{im}}\simeq M_2(K)
       \lhook\joinrel\longrightarrow
       \operatorname{End}^0_{\mathbb F_Q}(\mathcal A^2)=M_2(L),
                                                               \tag{P-F70c}
\]
and the actual source generators have the explicit images
\[
 Q_1\longmapsto
       \begin{pmatrix}f&\gamma\\0&0\end{pmatrix},\qquad
 Q_2\longmapsto
       \begin{pmatrix}0&0\\1&f\end{pmatrix}.                    \tag{P-F70d}
\]
In particular the actual prescribed central word satisfies
\[
                  FZ\longmapsto a I_2
                  =\frac{\Pi+\mathsf V}{2q^{26}} I_2.          \tag{P-F70e}
\]
The images of the rescaled generators $FQ_i$ are $F$ times
the matrices in (P-F70d), with no change of source normalization.

With the explicit opposite-algebra convention for pullback
described below, this is an all-word intertwiner on
$H^1(\mathcal A^2_{\overline{\mathbb F}_Q},\mathbb Q_\ell)$
for every $\ell\ne3$.  After $\ell$-adic scalar extension,
that eight-dimensional space is the
scalar extension of the actual $K$-module $V_q\simeq K^2$
from $K$ to $L$, and, after choosing a $K$-basis of $L$,
consists of two copies of the actual four-dimensional
rational source module.

For any polarization of $\mathcal A$ defined over $\mathbb F_Q$
and its product polarization on $\mathcal A^2$, the central
quasi-endomorphism $aI_2$ is Rosati self-adjoint and positive
in the following precise sense:
\[
 \operatorname{Tr}_{L/\mathbb Q}
    \operatorname{tr}_2(a X^\dagger X)>0
                      \quad\text{for every }0\ne X\in M_2(L).
                                                               \tag{P-F70f}
\]
Here $\dagger$ is the Rosati involution of that product
polarization and $\operatorname{tr}_2$ is the matrix trace.
\end{proposition}

\begin{proof}
\par\medskip\noindent\textbf{The rational source specialization and its full image.}
The lattice in (P-F52ax) has free basis
\[
 z_1=v_0,\quad z_2=\widehat v v_0,\quad
 z_3=Q_2v_0,\quad z_4=\widehat v Q_2v_0
\]
over $\mathbb Z[q,q^{-1}]$.  Its matrices for $FQ_1,FQ_2$
are the exact source matrices (P-F52ba), whose entries are
integer Laurent polynomials.  Specialization followed by
extension to $\mathbb Q$ therefore gives a four-dimensional
space with those specialized matrices.  Since $F(q)>0$,
division by $F(q)$ also defines its unscaled $Q_i$ actions.
This uses the actual source lattice and all its specialized
relations; it does not assume that a generic tensor-image
embedding remains faithful at every parameter.

On this module put $G=f^2I-(Q_1-Q_2)^2$.  P-F52 establishes
the exact identities
\[
 G^2-L_0G+N_0I=0,\qquad
 Z=S-G,\qquad FG=h+\widehat v,                                \tag{P-F70g}
\]
where $L_0,N_0,h$ are the source Laurent polynomials in
(P-F52f1), (P-F52aw); $L_0$ is the polynomial denoted $L$
in (P-F52f1), distinguished here from the CM field $L$.
The discriminant is
$L_0^2-4N_0=D/q^{40}$.  At $q=3^e$ this is nonsquare
over $\mathbb Q$ by P-F66, so the quadratic polynomial
is irreducible.  Hence $\mathbb Q[G]$ is a quadratic
field and $V_q$ is two-dimensional over it.

The source relation between its roots and $FZ$ identifies
this field with the $K$ in P-F69 by
\[
                      G\longmapsto \gamma=S-a/F.
                                                               \tag{P-F70h}
\]
Indeed $a$ has conjugates $F(S-g_1),F(S-g_2)$ by the
actual-word identity in P-F52, while P-F69 identifies
$\Pi+Q\Pi^{-1}$ with the element having conjugates
$2q^{26}\alpha_\pm(q)$.  All coefficients and signs in
(P-F70h) are therefore fixed by the source normalization.
The element $\gamma$ generates $K$, and
\[
                    \gamma\notin\{0,f^2\},
                                                               \tag{P-F70i}
\]
because neither of its distinct quadratic conjugates is rational.

The vectors $v_0,Q_2v_0$ form a $K$-basis of $V_q$.
To check independence directly, $Q_1v_0=fv_0$ and
$Q_1Q_2v_0=Gv_0$ by the source identities.  If
$Q_2v_0=t v_0$ with $t\in K$, then $Q_2^2=fQ_2$
gives $t\in\{0,f\}$ and the latter identity gives
$\gamma=ft\in\{0,f^2\}$, contradicting (P-F70i).
The resulting actual matrices are exactly (P-F70d).
Equivalently, the explicit map from the original four
source basis vectors to $K^2$ is
\[
 \begin{aligned}
 z_1&\longmapsto(1,0),&
 z_2&\longmapsto(F\gamma-h,0),\\
 z_3&\longmapsto(0,1),&
 z_4&\longmapsto(0,F\gamma-h).
 \end{aligned}                                                 \tag{P-F70j}
\]
Since $\gamma\notin\mathbb Q$ and $F\ne0$, this is an
invertible rational map, with the stated $K$-linear
interpretation.  The identity $FG=h+\widehat v$ checks
its normalization.

Write the matrices in (P-F70d) as $A_1,A_2$.
Their rational image algebra already contains $K$,
because $f^2I-(A_1-A_2)^2=\gamma I$.  The identities
\[
 \begin{split}
 fA_1-A_1A_2&=(f^2-\gamma)E_{11},\\
 fA_1A_2-\gamma A_1&=\gamma(f^2-\gamma)E_{12},\\
 E_{22}&=I-E_{11},\qquad E_{21}=A_2-fE_{22}
 \end{split}                                                   \tag{P-F70k}
\]
give all matrix units.  Their denominators are invertible
by (P-F70i).  Thus the image is all of $M_2(K)$.
This proves splitting and faithfulness on the stated
image algebra directly, rather than assuming that a
quadratic center forces a matrix algebra.

\par\medskip\noindent\textbf{The geometric algebra and the cohomological convention.}
P-F69 gives $L=\operatorname{End}^0_{\mathbb F_Q}(\mathcal A)$
and identifies $K\subset L$ through
$\beta=\Pi+Q\Pi^{-1}$.  The latter expression belongs to
the actual quasi-endomorphism algebra, since Frobenius
is an isogeny.  The standard product inclusions and
projections give
$\operatorname{End}^0_{\mathbb F_Q}(\mathcal A^2)=M_2(L)$.
Apply the inclusion $K\subset L$ to every matrix entry.
This proves the faithful algebra map $\jmath$ in (P-F70c).
Substituting $FZ=F(S-G)$ gives (P-F70e).

Pullback on $H^1$ is additive on homomorphisms and reverses
composition.  The canonical identification with the dual
Tate module gives this assertion also for rational
quasi-endomorphisms: Milne, \emph{Abelian Varieties},
Theorem~15.1(a), p.~27, in the version cited below.
To retain the
specified left source action and word order, use
\[
          \rho_\ell(X)=\bigl(\jmath(X^{\mathsf T})\bigr)^*,
                           \qquad X\in M_2(K).                \tag{P-F70l}
\]
Transposition is an anti-homomorphism here because $K$
is commutative.  Thus $\jmath(X^{\mathsf T})$ is the
explicit map into the opposite quasi-endomorphism
algebra, and contravariant pullback turns (P-F70l)
into a homomorphism:
\[
                         \rho_\ell(XY)=\rho_\ell(X)\rho_\ell(Y).
\]
For a matrix of maps on a product, pullback transposes
the component matrix.  Consequently (P-F70l) acts by
the original matrix $X$ on
$H^1(\mathcal A)\oplus H^1(\mathcal A)$, with each
entry of $L$ acting by its pullback on $H^1(\mathcal A)$.
This also proves the equality for every ordered word,
including the source's mixed relation, at the level of
these actual cohomological operators.

Put $K_\ell=K\otimes\mathbb Q_\ell$ and
$L_\ell=L\otimes\mathbb Q_\ell$.  P-F69 shows that
$H^1(\mathcal A,\mathbb Q_\ell)$ is free of rank one over
$L_\ell$.  Choose one generator $\eta$ of that module.
Then the explicit source-module isomorphism is
\[
 \begin{split}
 (V_q\otimes_{\mathbb Q}\mathbb Q_\ell)
                   \otimes_{K_\ell}L_\ell
       &\longrightarrow H^1(\mathcal A^2,\mathbb Q_\ell),\\
 (x,y)\otimes t&\longmapsto(xt\eta,yt\eta),
 \end{split}                                                   \tag{P-F70m}
\]
where (P-F70j) identifies the source pair $(x,y)$.
It intertwines $Q_i$ with $\rho_\ell(A_i)$, including
their factors $F$, and hence all source words and sums.
Since $L$ is quadratic over $K$, choosing the basis
$1,\Pi$ gives two copies of $V_q\otimes\mathbb Q_\ell$.
This does not identify $V_q$ with the whole
eight-dimensional target without that multiplicity.

Already on $H^1(\mathcal A)$, multiplication by $a$ is
two copies of the quadratic central multiplication,
so it realizes the actual central $FZ$ action on
$V_q$ after scalar extension and a choice of $K_\ell$-basis.
The passage to $\mathcal A^2$ is what realizes both
noncommuting generators by geometric quasi-endomorphisms.

\par\medskip\noindent\textbf{Rosati positivity of the actual central correspondence.}
Choose an ample line bundle on $\mathcal A$ over $\mathbb F_Q$
and its polarization.  Such a choice exists because an
abelian variety is projective over its base field.
Milne, \emph{Abelian Varieties}, Theorem~7.1, p.~10,
and Lemma~19.2, p.~41, give projectivity and the identity
\[
                          \Pi^\dagger\Pi=Q.
\]
The author-maintained corrected version of the original
article retains its theorem numbering:
\url{https://www.jmilne.org/math/xnotes/AVs.pdf}.
Because $L=\mathbb Q(\Pi)$, its Rosati involution is
therefore exactly the CM conjugation
$\Pi\mapsto Q/\Pi$.  It fixes $K$ and $a$.
For the product polarization, the involution on
$M_2(L)$ is conjugate transpose.  Thus
\[
 \operatorname{Tr}_{L/\mathbb Q}
       \operatorname{tr}_2(a X^\dagger X)
   =\sum_{\sigma:L\hookrightarrow\mathbb C}
          \sigma(a)\sum_{i,j}|\sigma(X_{ij})|^2>0
                                                               \tag{P-F70n}
\]
for $X\ne0$: all real conjugates of $a$ are the positive
$\alpha_\pm$ proved in P-F39, and a nonzero field entry
has no zero embedding.  This is compatible with the
standard positive Rosati trace form of Milne,
Theorem~17.3, p.~35, and explicitly proves the stronger
weighted form required here.  It is a positivity statement
for this auxiliary correspondence, not a claim that it
is Frobenius or an effective algebraic cycle.
\end{proof}

\paragraph{What has and has not been identified.}
The map realizes the eight-dimensional rational image
$M_2(K)$ of the actual specialized four-dimensional source
module.  It is not asserted to be faithful on the entire
$21$-dimensional source algebra.  It preserves the actual
two generator matrices, their unit, ordered products,
relations and prescribed scalar normalization; the
opposite-algebra convention in (P-F70l) is part of the
cohomological comparison.

The ambient $H^1$ is pure of weight one, but the operators
$Q_i,FZ$ are rational quasi-endomorphisms commuting with
Frobenius, not Frobenius itself.  The scalar $a$ is
$(\Pi+\mathsf V)/(2q^{26})$, not a Tate-twisted Frobenius.
Consequently P-F66's eigenvalue obstruction does not apply
to this auxiliary action.  No nonnegative structure
constants for a source candidate canonical basis, compatible
graded coefficient objects, or effective correspondences
have been constructed.  The parameter $q$ has been
specialized; the construction does not geometrize its
formal grading.  The base-extension incompatibility in
P-F69 remains, and no uniform family or comparison with
the supplied iterated theta operations follows.
